\section{Introduction}
\IEEEPARstart{T}{he} decision to enrol at a particular engineering college
is made, in large part, before the applicant ever sets foot on the campus.
Prospective students and their families gather information from institutional
websites, PDF brochures, ranking portals, and word of mouth, and they form
a mental model of the campus --- its departments, laboratories, hostel,
library, transport, and physical layout --- from fragmentary and often
inconsistent sources. Two information needs dominate this process. The
first is \emph{propositional}: what are the admission routes, the eligibility
criteria, the fee structure, the programs on offer, the accreditation
status? The second is \emph{spatial}: where is the Computer Science block,
which floor is a given laboratory on, what does the central quadrangle
actually look like, and how would one walk from the main gate to the
auditorium?

Digital campus experiences typically serve one of these needs well and the
other poorly. Institutional websites and chatbots answer propositional
questions but present the campus as a list of building names. Virtual-tour
products render 360$^{\circ}$ imagery but carry no queryable knowledge and
no notion of a route. Interactive campus maps show geography but not the
institution's policies or its interiors. A visitor who wants to know
\emph{``which floor is the Power System Simulation Lab on, and can I see
it?''} is served by neither category alone.

At the same time, the natural technology for a unified conversational
interface --- a large language model (LLM) --- introduces a well-documented
failure mode. LLMs asked institution-specific questions readily produce
fluent but fabricated answers: invented faculty names, wrong phone numbers,
plausible-sounding but incorrect fee figures~\cite{ji2023hallucination}.
For a system whose entire value proposition is being an authoritative
substitute for the official sources, an ungrounded answer is worse than no
answer. Retrieval-augmented generation (RAG)~\cite{lewis2020rag} mitigates
this by conditioning the model on retrieved evidence, but a naive RAG
pipeline still hallucinates when retrieval is weak, when the evidence is
thin, or when the model over-extrapolates from a partially relevant
passage~\cite{shuster2021retrieval,gao2023ragsurvey}.

\subsection{The \system{} System}
This paper describes \system{}, an integrated platform developed for the
Global Academy of Technology (GAT), Bengaluru, that combines the
propositional and spatial modalities in a single web application and treats
groundedness as a first-class engineering requirement. The system comprises
three cooperating subsystems:

\begin{enumerate}
\item A \textbf{grounded RAG assistant} that answers natural-language
questions using a knowledge base built exclusively from official
institutional sources ($128$ PDFs and $42$ web pages, $1{,}507$ text
chunks). Retrieval is hybrid (dense embeddings plus BM25); generation is
performed by a locally hosted Llama~3.2 model constrained by a
grounding-only system prompt, a confidence gate that can refuse before the
model is ever called, and a post-generation verification pass that
withholds answers containing unverifiable numeric or contact details.
\item A \textbf{360$^{\circ}$ virtual tour} of the institution's main
building, structured as a per-floor doubly linked list of panorama scenes
with graph-derived branch connections and a construction pass that stitches
floor boundaries so that both manual and automatic (``guided'') traversal
cross between floors with no floor-boundary-specific code.
\item A \textbf{campus navigation engine} implementing A$^{\star}$ shortest
paths over a $180$-node, $325$-edge graph that is rebuilt on demand from
the live PostgreSQL relational schema, with a heuristic that adapts to the
coordinate information actually available for each node pair.
\end{enumerate}

Query routing between the RAG assistant's specialist agents is performed by
a deterministic phrase- and keyword-matching supervisor, with a small
trained LSTM intent classifier consulted only as a fallback for phrasings
the deterministic rules do not cover. Voice input is provided through the
browser-native Web Speech API and feeds the same text pipeline as typed
input.

\subsection{Contributions}
This work makes the following contributions, stated conservatively and
tied directly to the implemented artifact:

\begin{enumerate}
\item A \textbf{system-level architecture} that unifies grounded
conversational question answering, 360$^{\circ}$ scene-graph tour
navigation, and graph-based wayfinding for a single institution, with all
three subsystems sharing one relational campus model and one web/service
stack (Section~\ref{sec:architecture}).
\item A detailed, \textbf{code-verified algorithmic description} of a
layered hallucination-mitigation design for institutional RAG: query
expansion for candidate generation, weighted dense--lexical score fusion,
a deterministic reranker, a two-term confidence gate, a grounding-only
prompt, and an independent post-generation numeric-claim check
(Section~\ref{sec:rag}, Algorithms~\ref{alg:kb}--\ref{alg:grounding}).
\item A \textbf{deterministic-primary, ML-fallback multi-agent routing}
scheme with full decision logging, and an honest account of a trained
LSTM intent classifier integrated as a subordinate signal rather than the
primary router (Section~\ref{sec:agents},
Algorithms~\ref{alg:routing}--\ref{alg:lstm}).
\item A \textbf{doubly-linked-list panorama scene graph} with a
tail-to-head cross-floor stitching mechanism that renders floor-boundary
traversal transparent to every consumer, and a six-phase guided-tour
traversal automaton with cancellation-token and cycle-guard safety
(Section~\ref{sec:tour}, Algorithms~\ref{alg:tourwalk}--\ref{alg:guided}).
\item A \textbf{database-backed A$^{\star}$ wayfinding engine} with a
three-tier degrading heuristic, together with a graph-construction
procedure that treats a normalized relational schema as the live source of
truth for the navigation graph (Section~\ref{sec:nav},
Algorithms~\ref{alg:graphbuild}--\ref{alg:astar}).
\item A set of \textbf{verified implementation-scale measurements} and a
\textbf{proposed evaluation protocol} (Section~\ref{sec:eval}) covering
retrieval quality, reranker ablation, intent-classifier generalization,
answer faithfulness, confidence-threshold calibration, pathfinding
performance, and usability --- experiments that the current artifact does
not yet include and that this paper does not pre-empt with fabricated
numbers.
\end{enumerate}

\subsection{Scope and Honesty Statement}
Consistent with reproducible-systems practice, we distinguish throughout
between (a)~what is implemented, (b)~what has been measured, (c)~what is
proposed for evaluation, and (d)~what is future work. Several components are
partial: the support-vector reranker is implemented but never trained and
is not the production path; the LSTM classifier's probability is not yet
wired into the confidence estimator; the A$^{\star}$ engine has no
front-end consumer; and a previously built three-dimensional map and
GPS turn-by-turn interface were deliberately removed from the codebase.
These facts are reported where relevant rather than elided.

The remainder of the paper is organized as follows.
Section~\ref{sec:related} surveys related work across ten areas and
positions \system{} against it. Section~\ref{sec:problem} formalizes the
problem and objectives. Section~\ref{sec:architecture} presents the system
architecture. Sections~\ref{sec:dataset}--\ref{sec:nav} detail the campus
knowledge representation, the virtual tour engine, the RAG pipeline, the
multi-agent and intent-classification subsystem, and the navigation engine.
Section~\ref{sec:voice} covers voice interaction and
Section~\ref{sec:impl} the implementation and technology stack.
Section~\ref{sec:eval} defines the proposed experimental methodology;
Section~\ref{sec:results} reports the measurements that currently exist and
discusses them. Sections~\ref{sec:limitations}--\ref{sec:conclusion} give
limitations, future work, and conclusions.
